%% ====================================================================
%%  Plantilla de evaluación
%% ====================================================================
\section{Plantilla de evaluación}\label{sec:plantilla}

A continuación, se presentan las plantillas de evaluación NPO
construidas para el análisis de las tecnologías proxy-caché y de
almacenamiento de archivos de gran tamaño. Las calificaciones siguen una
escala de 1 a 3, donde 1 es ``Muy limitado'', 2 es ``Aceptable'' y 3 es
``Excelente / Líder''.

En esta plantilla se presentan 6 indicadores que tienen un valor
ponderado con respecto a la puntuación final de cada una de las
herramientas; los indicadores que se tienen son evaluados con respecto a
las métricas que se habían definido en el formulario para la
caracterización de la organización. Cada una de estas métricas se
definen en la sección~\ref{sec:criterios}, criterios de evaluación, en donde se detalla la
importancia de cada una de estas en el contexto del proyecto.

A continuación, se le da un sentido a cada uno de los indicadores con
respecto a las métricas.

\begin{itemize}
\item \textbf{Rendimiento técnico (25\%):} Repositorio local de archivos grandes, priorización de tráfico académico, tiempos de descarga deficientes, optimización de ancho de banda, caché de archivos grandes, volúmenes de transferencia altos y tiempos de respuesta promedio.
\item \textbf{Escalabilidad (20\%):} Repositorio local de archivos grandes, priorización de tráfico académico (QoS), uso no misional de red, gobernanza del tráfico, QoS y priorización, ACL listas blancas/negras, escalabilidad de almacenamiento, políticas formales y diferenciar tráfico misional/no misional.
\item \textbf{Soporte y Mantenimiento (20\%):} Automatización de certificados SSL, falta de certificados SSL firmados, gestión de certificados digitales y adaptación a cambios.
\item \textbf{Adopción y Madurez (10\%):} Automatización de certificados SSL, proxy transparente/inverso/estándar e integración con servicios existentes.
\item \textbf{Observabilidad (10\%):} Uso no misional de red, métricas de uso de red y proxy, toma de decisiones basada en datos, monitoreo en tiempo real y alta disponibilidad, consumo de ancho de banda por usuario, estadísticas de tráfico por categoría y reportes.
\item \textbf{Adecuación al Contexto GRID (10\%):} Ajuste operativo de la herramienta en el contexto de operaciones del GRID y su escalabilidad.
\item \textbf{Riesgo Tecnológico (5\%):} Falta de certificados SSL firmados y sección para los estados proyectados para el año 2026 en la Matriz DAR.
\end{itemize}

La Tabla~\ref{tab:plantilla-proxy} presenta la plantilla de análisis DAR empleada para evaluar
las herramientas de proxy-caché, con los criterios y sus respectivos
pesos.

\begin{table}[H]
\centering\footnotesize
\begin{tabular}{|>{\raggedright\arraybackslash\bfseries}p{2.4cm}|c|c|c|c|c|c|c|c|c|}
\hline
\rowcolor{gridhead} \thead[l]{Criterios/\\Herram.} & \thead{Frec.\\SMS} & \thead{Rend.\\(25\%)} & \thead{Esc.\\(20\%)} & \thead{Sop.\\(20\%)} & \thead{Adop.\\(10\%)} & \thead{Obs.\\(10\%)} & \thead{Adec.\\(10\%)} & \thead{Riesgo\\(5\%)} & \thead{DAR} \\ \hline
HAProxy & 0 &  &  &  &  &  &  &  &  \\ \hline
\rowcolor{gridrow} NGINX & 1-2 &  &  &  &  &  &  &  &  \\ \hline
Envoy & 1-2 &  &  &  &  &  &  &  &  \\ \hline
\rowcolor{gridrow} Varnish & 0 &  &  &  &  &  &  &  &  \\ \hline
Traefik & 0 &  &  &  &  &  &  &  &  \\ \hline
\rowcolor{gridrow} Caddy & 0 &  &  &  &  &  &  &  &  \\ \hline
Apache Traffic Server & 0 &  &  &  &  &  &  &  &  \\ \hline
\rowcolor{gridrow} Kong & 1 &  &  &  &  &  &  &  &  \\ \hline
Squid & 0 &  &  &  &  &  &  &  &  \\ \hline
\rowcolor{gridrow} Istio & 0 &  &  &  &  &  &  &  &  \\ \hline
Linkerd & 0 &  &  &  &  &  &  &  &  \\ \hline
\rowcolor{gridrow} K8Sidecar & 1 &  &  &  &  &  &  &  &  \\ \hline
Proxy-Service & 1 &  &  &  &  &  &  &  &  \\ \hline
\rowcolor{gridrow} nuster & 0 &  &  &  &  &  &  &  &  \\ \hline
\end{tabular}
\caption{Plantilla análisis DAR proxy-caché}
\label{tab:plantilla-proxy}
\end{table}

De igual manera, en la Tabla~\ref{tab:plantilla-almac} se muestra la plantilla de análisis DAR
aplicada a las herramientas de almacenamiento.

\begin{table}[H]
\centering\footnotesize
\begin{tabular}{|>{\raggedright\arraybackslash\bfseries}p{2.7cm}|c|c|c|c|c|c|c|c|c|}
\hline
\rowcolor{gridhead} \thead[l]{Herramienta} & \thead{Frec.\\SMS} & \thead{Rend.\\(25\%)} & \thead{Esc.\\(20\%)} & \thead{Sop.\\(20\%)} & \thead{Adop.\\(10\%)} & \thead{Obs.\\(10\%)} & \thead{Adec.\\(10\%)} & \thead{Riesgo\\(5\%)} & \thead{DAR} \\ \hline
Ceph & 0 &  &  &  &  &  &  &  &  \\ \hline
\rowcolor{gridrow} Lustre & 0 &  &  &  &  &  &  &  &  \\ \hline
JuiceFS & 0 &  &  &  &  &  &  &  &  \\ \hline
\rowcolor{gridrow} DAOS & 0 &  &  &  &  &  &  &  &  \\ \hline
SeaweedFS & 0 &  &  &  &  &  &  &  &  \\ \hline
\rowcolor{gridrow} ROOT / RDataFrame / Spark / Dask / OSCAR & 1 &  &  &  &  &  &  &  &  \\ \hline
HDFS (EC + ISA-L / SYCL) & 1 &  &  &  &  &  &  &  &  \\ \hline
\rowcolor{gridrow} Knative + x264 / AV1 & 1 &  &  &  &  &  &  &  &  \\ \hline
Longhorn & 0 &  &  &  &  &  &  &  &  \\ \hline
\rowcolor{gridrow} ExaHDF5 & 1 &  &  &  &  &  &  &  &  \\ \hline
HDF5 & 1-2 &  &  &  &  &  &  &  &  \\ \hline
\rowcolor{gridrow} MinIO & 0 &  &  &  &  &  &  &  &  \\ \hline
HDFS & 1 &  &  &  &  &  &  &  &  \\ \hline
\rowcolor{gridrow} RustFS & 0 &  &  &  &  &  &  &  &  \\ \hline
BeeOND & 1 &  &  &  &  &  &  &  &  \\ \hline
\rowcolor{gridrow} OpenStack Swift & 0 &  &  &  &  &  &  &  &  \\ \hline
Garage & 0 &  &  &  &  &  &  &  &  \\ \hline
\rowcolor{gridrow} IPFS & 1 &  &  &  &  &  &  &  &  \\ \hline
GekkoFS & 1 &  &  &  &  &  &  &  &  \\ \hline
\rowcolor{gridrow} GlusterFS & 0 &  &  &  &  &  &  &  &  \\ \hline
BurstFS & 1 &  &  &  &  &  &  &  &  \\ \hline
\rowcolor{gridrow} TEES & 1 &  &  &  &  &  &  &  &  \\ \hline
\end{tabular}
\caption{Plantilla análisis DAR almacenamiento}
\label{tab:plantilla-almac}
\end{table}
